\documentclass[11pt,a4paper]{article}
\usepackage[utf8]{inputenc}
\usepackage[T1]{fontenc}
\usepackage{amsmath,amssymb,amsthm}
%\usepackage{stmaryrd}
\usepackage{geometry}
\usepackage{enumitem}
\usepackage{booktabs}
\usepackage{listings}
\usepackage{xcolor}
\usepackage{hyperref}

\geometry{margin=2.5cm}

\lstdefinestyle{racer}{
  basicstyle=\ttfamily\small,
  keywordstyle=\bfseries,
  commentstyle=\color{gray}\itshape,
  breaklines=true,
  frame=single,
  backgroundcolor=\color{gray!5},
  xleftmargin=1em,
  framexleftmargin=0.5em,
  columns=fullflexible,
  literate={->}{$\rightarrow$}1
}

\newcommand{\DLsome}[2]{\exists #1.\,#2}
\newcommand{\DLall}[2]{\forall #1.\,#2}
\newcommand{\DLand}{\sqcap}
\newcommand{\DLor}{\sqcup}
\newcommand{\DLnot}{\neg}
\newcommand{\DLsub}{\sqsubseteq}
\newcommand{\DLbot}{\bot}
\newcommand{\DLtop}{\top}
\newcommand{\Racer}{\textsc{Racer}}
\newcommand{\concept}[1]{\ensuremath{\mathsf{#1}}}
\newcommand{\role}[1]{\ensuremath{\mathit{#1}}}

\title{Using Standard Description Logic Reasoning\\for SAT Planning:\\
Exploring Role-Based Encodings of the\\Wolf--Goat--Cabbage Riddle}

\author{Michael Wessel \and Claude (Anthropic)}

\date{March 2026}

\begin{document}

\maketitle

\begin{abstract}
SAT planning reduces planning problems to satisfiability tests: a plan
of length~$n$ exists if and only if the conjunction of start state, goal
state, and action axioms is satisfiable. The classical propositional
encoding suffers from \emph{symbol proliferation}---every combination of
object and location requires its own propositional variable. We
investigate whether Description Logic (DL) roles can yield a more
compact encoding, using the well-known Wolf--Goat--Cabbage river
crossing puzzle as a case study. Starting from an existing encoding with
atomic concept names in KRSS syntax for the \Racer{} reasoner, we
systematically explore several role-based reformulations, documenting
the trade-offs between symbol economy, reasoning performance, and
semantic correctness.  Our journey reveals fundamental tensions between
compactness and tableau efficiency in DL-based SAT planning.
\end{abstract}

%% ============================================================
\section{Introduction}
\label{sec:intro}
%% ============================================================

The Wolf--Goat--Cabbage (WGC) riddle is a classical river crossing
puzzle: a ferryman must transport a wolf, a goat, and a cabbage across a
river.  The boat holds at most the ferryman plus one object.  The wolf
and goat cannot be left alone (the wolf eats the goat), nor can the goat
and cabbage (the goat eats the cabbage).

Following the SAT planning paradigm~\cite{kautz1992}, the problem can
be encoded as a satisfiability test.  States are connected by a
functional \role{next} relation; axioms constrain which successor states
are legal.  A plan of length~$n$ exists iff the conjunction of axioms,
start state, and goal state is satisfiable for a chain of $n$~states.

Wessel~\cite{wessel2014} presented an encoding using the \Racer{}
Description Logic reasoner, employing atomic concept names in KRSS
syntax.  While effective, this encoding uses over 20 atomic concept names
for object--location combinations (e.g., \concept{goat\text{-}a},
\concept{wolf\text{-}boat\text{-}from\text{-}a\text{-}to\text{-}b}),
exhibiting the symbol proliferation typical of propositional SAT
encodings.

In this paper, we explore whether DL \emph{roles} can provide a more
compact representation, and document the challenges encountered along
the way.

%% ============================================================
\section{The Original Encoding}
\label{sec:original}
%% ============================================================

\subsection{Logical Signature}

Each object (goat, wolf, cabbage, ferryman, boat) can be in one of four
locations: riverbank~$A$, riverbank~$B$, in the boat going from~$A$
to~$B$, or in the boat going from~$B$ to~$A$.  This yields four atomic
concept names per object, for example:
\[
\concept{goat\text{-}a}, \quad
\concept{goat\text{-}b}, \quad
\concept{goat\text{-}boat\text{-}from\text{-}a\text{-}to\text{-}b}, \quad
\concept{goat\text{-}boat\text{-}from\text{-}b\text{-}to\text{-}a}
\]
The role \role{next} is declared as a \emph{feature} (functional
role), connecting each state to its unique successor.

\subsection{Successor State Axioms}

The boat cycles deterministically through the four locations:
\[
\concept{boat\text{-}a} \DLsub \DLall{\role{next}}{\concept{boat\text{-}from\text{-}a\text{-}to\text{-}b}}
\]
\[
\concept{boat\text{-}from\text{-}a\text{-}to\text{-}b} \DLsub \DLall{\role{next}}{\concept{boat\text{-}b}}
\]
and so on.  In \Racer{} KRSS syntax:
\begin{lstlisting}[style=racer]
(implies boat-a (all next boat-from-a-to-b))
(implies boat-from-a-to-b (all next boat-b))
(implies boat-b (all next boat-from-b-to-a))
(implies boat-from-b-to-a (all next boat-a))
\end{lstlisting}

\subsection{No-Teleportation Axioms}

Objects can only move by boarding the boat. From a riverbank, an object
either stays or boards:
\[
\concept{goat\text{-}a} \DLsub
  \DLall{\role{next}}{\concept{goat\text{-}a}} \DLor
  \DLall{\role{next}}{\concept{goat\text{-}boat\text{-}from\text{-}a\text{-}to\text{-}b}}
\]
From transit, it arrives at the opposite bank:
\[
\concept{goat\text{-}boat\text{-}from\text{-}a\text{-}to\text{-}b} \DLsub
  \DLall{\role{next}}{\concept{goat\text{-}b}}
\]
The same pattern is repeated for wolf, cabbage, and ferryman (16~axioms total).

\subsection{Boat Capacity, Safety, and Integrity}

\textbf{Capacity:} When the boat travels from $A$ to~$B$, the ferryman
is aboard, exactly one of \{wolf, goat, cabbage\} is aboard, and no two
passengers ride together.  When returning from $B$ to~$A$, the ferryman
may return alone.

\textbf{Safety:} Wolf and goat, or cabbage and goat, on the same
riverbank require the ferryman's presence:
\[
\concept{wolf\text{-}a} \DLand \concept{goat\text{-}a} \DLsub \concept{ferryman\text{-}a}
\]

\textbf{Disjointness:} Each object is at exactly one location:
\begin{lstlisting}[style=racer]
(disjoint goat-a goat-b goat-boat-from-a-to-b goat-boat-from-b-to-a)
\end{lstlisting}

\textbf{Persistence:} Objects do not disappear---every state specifies
each object's location:
\[
\concept{goat} \DLsub
  \concept{goat\text{-}a} \DLor \concept{goat\text{-}b} \DLor
  \concept{goat\text{-}boat\text{-}from\text{-}a\text{-}to\text{-}b} \DLor
  \concept{goat\text{-}boat\text{-}from\text{-}b\text{-}to\text{-}a}
\]

\textbf{Coupling:} An object on the boat implies the boat is in transit:
\[
\concept{goat\text{-}boat\text{-}from\text{-}a\text{-}to\text{-}b} \DLsub
  \concept{goat} \DLand \concept{boat\text{-}from\text{-}a\text{-}to\text{-}b}
\]

\subsection{States, Start, and Goal}

Every state either is a goal or has a successor, and specifies all objects:
\[
\concept{state} \DLsub
  (\concept{goal} \DLor \DLsome{\role{next}}{\concept{state}}) \DLand
  \concept{goat} \DLand \concept{wolf} \DLand \concept{cabbage} \DLand
  \concept{ferryman} \DLand \concept{boat}
\]
\[
\concept{start} \DLsub \concept{state} \DLand \concept{boat\text{-}a} \DLand
  \concept{goat\text{-}a} \DLand \concept{wolf\text{-}a} \DLand
  \concept{cabbage\text{-}a} \DLand \concept{ferryman\text{-}a}
\]
\[
\concept{goal} \DLsub \concept{state} \DLand \concept{boat\text{-}b} \DLand
  \concept{goat\text{-}b} \DLand \concept{wolf\text{-}b} \DLand
  \concept{cabbage\text{-}b} \DLand \concept{ferryman\text{-}b}
\]

\subsection{Performance Characteristics}

This encoding uses exclusively \emph{atomic concept names} on the
left-hand side of GCIs (General Concept Inclusions), enabling effective
\emph{absorption} by the tableau reasoner.  Absorbed GCIs are handled as
lazy unfolding rules rather than global constraints checked against every
individual---a significant performance advantage.

The encoding works well in \Racer{}: ABox consistency checking completes
quickly, and individual type retrieval yields the expected two solutions
to the puzzle (corresponding to transporting the cabbage or the wolf
after the initial goat crossing).

The price is symbol proliferation: $5 \text{ objects} \times 4 \text{
locations} = 20$ atomic concept names, plus 5 generic object concepts
and supporting names.

%% ============================================================
\section{Attempt~1: Object-Centric Location Roles}
\label{sec:attempt1}
%% ============================================================

\subsection{Idea}

Replace the 20 atomic object--location concepts with 5 \emph{functional
roles} (features) mapping states to locations, plus 4 shared location
concepts:
\begin{align*}
\concept{goat\text{-}a} &\;\leadsto\; \DLsome{\role{goat\text{-}loc}}{\concept{a}} \\
\concept{goat\text{-}b} &\;\leadsto\; \DLsome{\role{goat\text{-}loc}}{\concept{b}} \\
\concept{goat\text{-}boat\text{-}from\text{-}a\text{-}to\text{-}b} &\;\leadsto\; \DLsome{\role{goat\text{-}loc}}{\concept{from\text{-}a\text{-}to\text{-}b}}
\end{align*}

The primitive vocabulary shrinks from 20+ symbols to 9 (5~roles +
4~location concepts).  Disjointness of locations combined with
functional roles gives per-object location uniqueness for free:

\begin{lstlisting}[style=racer]
(define-primitive-role goat-loc :feature t)
(disjoint a b from-a-to-b from-b-to-a)
\end{lstlisting}

\subsection{Axiom Rewriting}

All axioms are rewritten using role expressions.  For example, the boat
cycle:
\[
\DLsome{\role{boat\text{-}loc}}{\concept{a}} \DLsub
  \DLall{\role{next}}{\DLsome{\role{boat\text{-}loc}}{\concept{from\text{-}a\text{-}to\text{-}b}}}
\]
Coupling becomes:
\[
\DLsome{\role{goat\text{-}loc}}{\concept{from\text{-}a\text{-}to\text{-}b}} \DLsub
  \DLsome{\role{boat\text{-}loc}}{\concept{from\text{-}a\text{-}to\text{-}b}}
\]

\subsection{Performance Problem}

\textbf{This encoding was too slow for practical use.}  The root cause:
every axiom with $\DLsome{\role{r}}{C}$ on the left-hand side is a
\emph{General Concept Inclusion} (GCI) that cannot be absorbed into the
concept hierarchy.  The original encoding had atomic concept names on the
left---cheaply absorbed.  The role-based encoding has ${\sim}30$
non-absorbable GCIs, each requiring global checking against every
individual in the tableau.

Additionally, the \texttt{:feature} declarations on all 5 location roles
imposed expensive merge operations whenever the tableau created multiple
fillers for the same role.

%% ============================================================
\section{Attempt~2: Defined Concept Abbreviations}
\label{sec:attempt2}
%% ============================================================

\subsection{Idea}

Introduce named concepts as abbreviations for the role expressions:
\begin{lstlisting}[style=racer]
(define-concept goat-at-a (some goat-loc a))
(define-concept goat-at-b (some goat-loc b))
;; ... 20 definitions total
\end{lstlisting}
Then use these names in axioms: \texttt{(implies goat-at-a (all next
goat-at-a) ...)}. This gives the reasoner atomic-like concept names for
internal optimization while preserving the role-based semantics.

\subsection{Problem}

\textbf{This reintroduces symbol proliferation.}  We now have 20 defined
concept names---the same count as the original's 20 primitive concepts.
While they are \emph{defined from} the 9-symbol role vocabulary rather
than independently primitive, the practical effect is the same: 20
names that the user must manage.  The compactness goal is defeated.

%% ============================================================
\section{Attempt~3: Dropping Features}
\label{sec:attempt3}
%% ============================================================

\subsection{Idea}

Remove the \texttt{:feature} declarations from the location roles to
eliminate merge operations.  Use \texttt{(at-most 1 role)} in the state
axiom instead---scoped to state individuals rather than globally
declared:

\begin{lstlisting}[style=racer]
(define-primitive-role goat-loc)  ;; NOT a feature

(implies state (and ...
                    (at-most 1 goat-loc)
                    (at-most 1 wolf-loc)
                    ...))
\end{lstlisting}

\subsection{Problem}

\textbf{Still too slow.}  The \texttt{(at-most~1)} restriction is
semantically equivalent to a feature constraint---the reasoner must
still track cardinality and perform merge operations.  The GCI
bottleneck ($\DLsome{\role{r}}{C}$ on axiom left-hand sides) remains
unchanged.  The \texttt{at-most} restrictions were subsequently removed
entirely, relying on the axiom structure to naturally produce at most one
filler per role (the tableau is lazy and only creates fillers when forced
by existential restrictions).

%% ============================================================
\section{Attempt~4: Location-Centric Roles}
\label{sec:attempt4}
%% ============================================================

\subsection{Idea}

Invert the role perspective: instead of objects tracking their locations
(\role{goat-loc}, \role{wolf-loc}, \ldots), locations track which
objects are at them.  Use 4~roles (\role{at\text{-}a},
\role{at\text{-}b}, \role{at\text{-}ab}, \role{at\text{-}ba}) and
5~object concepts (\concept{goat}, \concept{wolf}, \concept{cabbage},
\concept{ferryman}, \concept{boat}):

\begin{align*}
\concept{goat\text{-}a} &\;\leadsto\; \DLsome{\role{at\text{-}a}}{\concept{goat}} \\
\concept{wolf\text{-}boat\text{-}from\text{-}a\text{-}to\text{-}b} &\;\leadsto\; \DLsome{\role{at\text{-}ab}}{\concept{wolf}}
\end{align*}

The symbol count is unchanged (9), but axioms that reference a single
location now use a single role throughout.  For example, the safety
constraint:
\[
\DLsome{\role{at\text{-}a}}{\concept{wolf}} \DLand
\DLsome{\role{at\text{-}a}}{\concept{goat}} \DLsub
\DLsome{\role{at\text{-}a}}{\concept{ferryman}}
\]
All three expressions reference the same role \role{at\text{-}a}, which
may benefit the reasoner's internal indexing.

\subsection{Open-World Semantics and \texttt{(all next ...)}}

A subtle but important issue arose with the state axiom.  The original
formulation used $\DLsome{\role{next}}{\concept{state}}$ to ensure
non-goal states have a successor:
\begin{lstlisting}[style=racer]
(implies state (and (or goal (some next state)) ...))
\end{lstlisting}

Under open-domain semantics, $\DLsome{\role{next}}{\concept{state}}$
does \emph{not} necessarily constrain the explicitly modeled successor
(linked via \texttt{(related s1 s2 next)}).  The reasoner may create a
fresh anonymous individual to satisfy the existential, leaving the
asserted successor unconstrained.

The fix: use $\DLall{\role{next}}{\cdot}$ to push object-persistence
constraints to the modeled successor.  Since \role{next} is a feature,
$\DLall{\role{next}}{C}$ applies to the unique successor.  Additionally,
declaring \texttt{:domain state :range state} on the \role{next} role
ensures that all individuals connected by \role{next} are classified as
states:
\begin{lstlisting}[style=racer]
(define-primitive-role next :feature t :domain state :range state)

(implies state (and (or goal (some next top))
                    (all next (or (some at-a goat) ...))
                    ...))
\end{lstlisting}

\subsection{The Uniqueness Problem}

\textbf{Without uniqueness enforcement, the encoding is unsound.}

The location-centric approach cannot use features on the \role{at-X}
roles because multiple objects legitimately share a location (e.g.,
wolf, cabbage, and ferryman can all be at riverbank~$A$).

Without uniqueness constraints, the safety axioms become toothless.
Consider state~$s_2$ where the boat departs from~$A$ with the ferryman.
If the goat stays at~$A$ alongside the wolf, safety requires
$\DLsome{\role{at\text{-}a}}{\concept{ferryman}}$.  Without uniqueness,
the reasoner simply creates an \emph{additional}
\role{at\text{-}a}~filler of type \concept{ferryman}---even though the
ferryman is already on the boat
($\DLsome{\role{at\text{-}ab}}{\concept{ferryman}}$).  The ferryman
appears at two locations simultaneously, trivially satisfying the safety
constraint.

In the original encoding, this was prevented by:
\begin{lstlisting}[style=racer]
(disjoint ferryman-a ferryman-b
          ferryman-boat-from-a-to-b ferryman-boat-from-b-to-a)
\end{lstlisting}

\subsection{Solution: Per-Object Location Exclusion Axioms}

Since per-object uniqueness cannot be achieved via features or
\texttt{disjoint} on atomic names (we don't have them), we add explicit
exclusion axioms.  For each object and each location, being at that
location excludes all others:
\[
\DLsome{\role{at\text{-}a}}{\concept{goat}} \DLsub
  \DLnot\DLsome{\role{at\text{-}b}}{\concept{goat}} \DLand
  \DLnot\DLsome{\role{at\text{-}ab}}{\concept{goat}} \DLand
  \DLnot\DLsome{\role{at\text{-}ba}}{\concept{goat}}
\]

In \Racer{} syntax:
\begin{lstlisting}[style=racer]
(implies (some at-a goat)
  (and (not (some at-b goat))
       (not (some at-ab goat))
       (not (some at-ba goat))))
\end{lstlisting}

This requires $5 \text{ objects} \times 4 \text{ locations} = 20$
exclusion axioms---effectively re-encoding the original's disjointness
declarations as GCIs.

\subsection{Readable Labels}

Since the role-based encoding produces types like
$\DLsome{\role{at\text{-}a}}{\concept{goat}}$ rather than atomic names,
ABox queries would not yield human-readable labels.  We add 20 labeling
axioms that tag states with familiar names:
\begin{lstlisting}[style=racer]
(implies (some at-a goat) goat-a)
(implies (some at-ab wolf) wolf-boat-from-a-to-b)
;; ... etc.
\end{lstlisting}

These are cheap GCIs: the right-hand side is an atomic concept name (no
fillers created, just tagging).

\subsection{Final Encoding}

The final, working encoding uses:
\begin{itemize}[nosep]
  \item 4 location roles: \role{at\text{-}a}, \role{at\text{-}b}, \role{at\text{-}ab}, \role{at\text{-}ba} (non-features)
  \item 5 object concepts: \concept{goat}, \concept{wolf}, \concept{cabbage}, \concept{ferryman}, \concept{boat} (pairwise disjoint)
  \item 1 functional role: \role{next} with domain/range \concept{state}
  \item 20 per-object exclusion axioms for uniqueness
  \item 20 labeling axioms for readable output
  \item Standard axioms: boat cycle (4), capacity (2), no-teleportation (16), coupling (8), safety (4), state/start/goal (3)
\end{itemize}

The encoding produces correct results and runs at reasonable speed in \Racer{}.

%% ============================================================
\section{Discussion}
\label{sec:discussion}
%% ============================================================

\subsection{Symbol Count Comparison}

Table~\ref{tab:comparison} summarizes the primitive vocabulary of each
approach.

\begin{table}[ht]
\centering
\begin{tabular}{lcccc}
\toprule
\textbf{Approach} & \textbf{Roles} & \textbf{Concepts} & \textbf{Total Primitives} & \textbf{GCIs} \\
\midrule
Original (atomic)        & 1 & 20+ & 21+ & 0 (all absorbed)\\
Object-centric roles     & 6 & 4   & 10  & ${\sim}30$ \\
Defined abbreviations    & 6 & 24  & 30  & ${\sim}30$ \\
Location-centric (final) & 5 & 5   & 10  & ${\sim}50$ \\
\bottomrule
\end{tabular}
\caption{Comparison of encoding approaches. GCI count refers to
non-absorbable axioms (complex concept on LHS).}
\label{tab:comparison}
\end{table}

\subsection{The Compactness--Performance Trade-off}

Our experiments reveal a fundamental tension:

\begin{enumerate}[nosep]
\item \textbf{Atomic concepts absorb; role expressions do not.}
  Axioms of the form $A \DLsub C$ where $A$ is atomic are absorbed
  into the concept hierarchy and handled as lazy unfolding.  Axioms
  of the form $\DLsome{\role{r}}{C} \DLsub D$ are GCIs checked
  globally against every tableau individual---far more expensive.

\item \textbf{Features are expensive but necessary.}
  Functional role declarations enable automatic per-object
  uniqueness but impose merge operations.  Dropping features
  requires explicit exclusion axioms (20 additional GCIs) to
  maintain soundness.

\item \textbf{Open-world semantics require care.}
  The use of $\DLsome{\role{next}}{\concept{state}}$ in the state
  axiom can lead to the creation of anonymous successor individuals
  rather than constraining the explicitly asserted ones.  Using
  $\DLall{\role{next}}{\cdot}$ with domain/range declarations
  ensures constraints apply to the modeled plan states.

\item \textbf{Location-centric roles cannot use features for
  uniqueness.}
  When multiple objects share a location role, features would
  incorrectly limit the role to a single filler.  Per-object
  uniqueness must be encoded separately.
\end{enumerate}

\subsection{Lessons Learned}

The original encoding, despite its symbol proliferation, is
well-matched to tableau-based DL reasoning.  Atomic concept names
enable absorption, \texttt{disjoint} declarations are handled
efficiently, and no GCIs are needed.

The role-based encodings achieve greater \emph{ontological economy}
(fewer primitive symbols), but at the cost of GCIs that the tableau must
check globally.  The location-centric variant adds locality benefits
(same-role references within safety and capacity axioms) and produces
correct results at reasonable speed, but requires 20~exclusion axioms
and 20~labeling axioms---totaling 40~additional axioms---that partially
offset the compactness gains.

The exercise demonstrates that in DL-based SAT planning, the choice of
representation is not merely aesthetic: it has direct consequences for
reasoning performance.  The ``right'' encoding depends on the reasoner's
optimization strategies, particularly GCI absorption and feature
handling.

%% ============================================================
\section{Conclusion}
\label{sec:conclusion}
%% ============================================================

We explored several role-based reformulations of a classical SAT
planning problem in Description Logics, seeking to reduce symbol
proliferation while maintaining reasoning efficiency.  Our journey
through object-centric roles, defined abbreviations, feature
elimination, and location-centric roles revealed that:

\begin{itemize}[nosep]
\item Compact role-based encodings are possible (9~primitive symbols
  vs.\ 20+), but produce non-absorbable GCIs that slow tableau
  reasoning.
\item Uniqueness of object locations---trivially handled by
  \texttt{disjoint} in the atomic encoding---requires either features
  (expensive), cardinality restrictions (equally expensive), or
  explicit exclusion axioms (verbose) in role-based encodings.
\item Open-world semantics demand careful use of
  $\DLall{\role{next}}{\cdot}$ rather than
  $\DLsome{\role{next}}{\cdot}$ to ensure constraints apply to
  explicitly modeled successor states.
\item The final location-centric encoding with exclusion axioms
  provides a working, reasonably fast alternative that demonstrates
  the use of DL roles for structuring planning state descriptions.
\end{itemize}

The original atomic encoding remains the pragmatic choice for
performance-critical applications.  However, the role-based approach
offers conceptual clarity and extensibility---for instance, adding new
objects or locations requires only new concept names, not a combinatorial
explosion of new atomic symbols.

\begin{thebibliography}{9}
\bibitem{kautz1992}
  H.~Kautz and B.~Selman,
  ``Planning as Satisfiability,''
  in \textit{Proc.\ ECAI}, 1992, pp.~359--363.

\bibitem{wessel2014}
  M.~Wessel,
  ``SAT Planning in Description Logics: Solving the Classical Wolf Goat Cabbage Riddle,''
  2014. Unpublished presentation.

\bibitem{haarslev2012}
  V.~Haarslev, K.~Hidde, R.~M{\"o}ller, and M.~Wessel,
  ``The \textsc{RacerPro} Knowledge Representation and Reasoning System,''
  \textit{Semantic Web -- Interoperability, Usability, Applicability},
  vol.~3, no.~3, pp.~267--277, IOS Press, 2012.
  \url{https://www.semantic-web-journal.net/content/racerpro-knowledge-representation-and-reasoning-system}
\end{thebibliography}

\end{document}
